%% kapa: exam=1 round=32 year=2021 period=1 form=A subject=civil_law subject_ko=민법 seq=31 number=31
%% source: Q-Net 감정평가사 기출문제 (cst00309) artlSeq=5208397 file=extracted/1008/5208397/A형/2021년도 제32회 감정평가사 1차 1교시 A형.pdf
%% conversion: kordoc 4.24.51 -> tools/gen1.py (auto); review=proofread-sonnet; flags=none
\begin{question}{31}
부동산 실권리자명의 등기에 관한 법률상 명의신탁에 관한 설명으로 옳은 것은? (다툼이 있으면 판례에 따름)
\begin{choices}
  \item 투기·탈세 등의 방지라는 법의 목적상 명의신탁은 그 자체로 선량한 풍속 기타 사회 질서에 위반된다.
  \item 명의신탁이 무효인 경우, 신탁자와 수탁자가 혼인하면 명의신탁약정이 체결된 때로부터 위 명의신탁은 유효하게 된다.
  \item 부동산 명의신탁약정의 무효는 수탁자로부터 그 부동산을 취득한 악의의 제3자에게 대항할 수 있다.
  \item 농지법에 따른 제한을 피하기 위하여 명의신탁을 한 경우에도 그에 따른 수탁자 명의의 소유권이전등기가 불법원인급여라고 할 수 없다.
  \item 조세포탈 등의 목적 없이 종교단체장의 명의로 그 종교단체 보유 부동산의 소유권을 등기한 경우, 그 단체와 단체장 간의 명의신탁약정은 유효하다.
\end{choices}
\end{question}
\answer{4}
